\centerline{\bf \Huge Финальная МатДрака}

\textit{ВСЕ участники команды, занявшей I место, получают по 8 баллов в рейтинговую таблицу и право на 8 попыток ввода кода.}

\textit{ВСЕ участники команды, занявшей II место, получают по 4 балла в рейтинговую таблицу и право на 4 попытки ввода кода.}

\textit{ВСЕ участники команды, занявшей III место, ничего не получают.}

\textbf{Ученик, решивший задачу, получает 6 баллов в рейтинговую таблицу и право на 6 попыток ввода кода; если решали двое, то по 3 балла и право на 3 попытки, если решали трое, то по 2 балла и право на 2 попытки}\\

\problem{\textbf{1}} На острове живут рыцари и лжецы. Рыцари всегда говорят правду,
а лжецы всегда лгут. Однажды 6 жителей острова собрались вместе и каждый сказал: «Среди
остальных пятерых ровно четыре лжеца!». Сколько рыцарей могло среди них быть?\\


\problem{\textbf{2}} На острове живут рыцари, которые всегда говорят правду, и
лжецы, которые всегда лгут. Однажды 15 аборигенов, среди которых были как рыцари, так и
лжецы, встали в хоровод, и каждый произнёс: «Из двух людей, стоящих напротив меня, один —
рыцарь, а другой — лжец». Сколько среди них рыцарей?
\\

\problem{\textbf{3}} Страна имеет форму квадрата и разделена на 25 одинаковых
квадратных графств. В каждом графстве правит либо граф-рыцарь, который всегда говорит
правду, либо граф-лжец, который всегда лжет. Однажды каждый граф сказал: «Среди моих
соседей поровну рыцарей и лжецов». Какое максимальное число рыцарей могло быть? (Графы
являются соседями, если их графства имеют общую сторону.)\\

\problem{\textbf{4}} Пять футбольных команд провели турнир — каждая команда сыграла с каждой по разу. За победу начислялось 3 очка, за ничью — 1 очко,
за проигрыш очков не давалось. Четыре команды набрали соответственно 1, 2, 5 и 7 очков.
А сколько очков набрала пятая команда?\\

\problem{\textbf{5}} Петя, Саша и Миша играют в теннис на вылет. Игра на вылет
означает, что в каждой партии играют двое, а третий ждёт. Проигравший партию уступает
место третьему и в следующей партии сам становится ждущим. Петя сыграл всего 12 партий,
Саша — 7 партий, Миша — 11 партий. Сколько раз Петя выиграл у Саши?\\




\problem{\textbf{6}} Сколько существует чисел, больших, чем 3528, каждое из которых можно получить перестановкой цифр данного числа?\\






\problem{\textbf{7}} Сколько существует четырёхзначных чисел, сумма цифр которых больше 33?\\



\problem{\textbf{8}} Килограмм говядины с костями стоит 78 рублей,
килограмм говядины без костей — 90 рублей, а килограмм костей — 15 рублей. Сколько граммов
костей в килограмме говядины?
\\


\problem{\textbf{9}}  Цена стандартного обеда в таверне «Буратино» зависит только от дня недели. Аня обедала 10 дней подряд, начиная с 10 июля, и
заплатила 70 сольдо. Ваня также заплатил 70 сольдо за 12 обедов, начиная с 12 июля. Таня
заплатила 100 сольдо за 20 обедов, начиная с 20 июля. Сколько заплатит Саня за 24 обеда,
начиная с 24 июля?
\\

\problem{\textbf{10}} Часы Безумного Шляпника спешат на 15 минут
в час, а часы Мартовского Зайца отстают на 10 минут в час. Однажды они поставили свои часы
по часам Сони (которые остановились и всегда показывают 12:00) и договорились собраться
в 5 часов вечера на традиционный файв-о-клок. Сколько времени Безумный Шляпник будет
ждать Мартовского Зайца, если каждый приходит ровно в 17:00 по своим часам?\\

\problem{\textbf{11}} Напишите вместо семи звёздочек семь различных цифр так, чтобы
получилось верное равенство: **** + ** + * = 2026.\\

\problem{\textbf{12}} Из последовательности натуральных
чисел 1, 2, 3, . . . удалили все точные квадраты ($1=1\times1,4=2\times2,9=3\times3...$). Какое число будет
находиться на 2026 месте среди оставшихся?






